\subsection{Chen \& Ong (2022): M3GNet universal graph IAP}
\label{sec:appendix-example-chen2022}

\paragraph{Q1 -- Bibliographic identification.}
Chen and Ong introduce M3GNet, a materials graph network with explicit
three-body interactions trained as a universal interatomic potential
across the periodic table. Published in \emph{Nat.\ Comput.\ Sci.} 2,
718--728 (2022); arXiv:2202.02450.
\lstinputlisting[language=turtle]{artifacts/kg/listings/chen2022/q01-bibliographic.ttl}

\paragraph{Q2 -- Material system.}
The MLIP covers Materials Project crystals across 89 elements
(62{,}783 distinct compounds spanning binaries, ternaries, and
quaternaries from the MPF.2021.2.8 subset). We model this as a single
$\MaterialSystemC$ at the database-coverage level rather than
enumerating individual compounds.
\lstinputlisting[language=turtle]{artifacts/kg/listings/chen2022/q02-material.ttl}

\paragraph{Q3 -- Reference calculation method.}
Reference data come from DFT calculations.
\lstinputlisting[language=turtle]{artifacts/kg/listings/chen2022/q03-reference-method-type.ttl}

\paragraph{Q4 -- Reference settings.}
DFT settings inherit the Materials Project defaults: VASP with PBE
GGA (or PBE+U for transition-metal-containing compounds) and PAW
pseudopotentials. The energy cutoff and k-mesh follow the per-task MP
workflow rather than a single fixed value.
\lstinputlisting[language=turtle]{artifacts/kg/listings/chen2022/q04-reference-settings.ttl}

\paragraph{Q5 -- Training dataset.}
MPF.2021.2.8 contains 187{,}687 ionic-step configurations from 62{,}783
compounds, covering energies, forces, and stresses. Provenance is
\emph{published} (the dataset is reused from MP).
\lstinputlisting[language=turtle]{artifacts/kg/listings/chen2022/q05-dataset.ttl}

\paragraph{Q6 -- Sampling strategies.}
Two strategies were combined: structural-relaxation trajectory
sampling (first and middle ionic steps of the first relaxation plus
the last step of the second relaxation, with high-energy or
short-distance frames filtered out), and periodic-table chemical-space
coverage (89 elements; 90/5/5 train/val/test split by material rather
than by configuration to avoid leakage).
\lstinputlisting[language=turtle]{artifacts/kg/listings/chen2022/q06-sampling.ttl}

\paragraph{Q7 -- MLIP method, components, implementation.}
M3GNet is a materials graph neural network with explicit three-body
angular interactions; energy is a sum of atomic contributions from a
gated MLP readout, with forces and stresses obtained by
auto-differentiation. Loss is Huber ($\delta=0.01$) on energy, forces,
and stresses with weights $w_E=1$, $w_f=1$, $w_\sigma=0.1$. Optimisation
is Adam with cosine learning-rate decay; the implementation is the
m3gnet TensorFlow package.
\lstinputlisting[language=turtle]{artifacts/kg/listings/chen2022/q07-method.ttl}

\paragraph{Q8 -- Hyperparameter settings.}
Reported settings: 2-body cutoff $r_c = 5.0$\,\AA, 3-body cutoff
$4.0$\,\AA, embedding dimension 64, 3 graph-network blocks, 3 radial
basis functions.
\lstinputlisting[language=turtle]{artifacts/kg/listings/chen2022/q08-hyperparameter-settings.ttl}

\paragraph{Q9 -- MLIP run and trained model.}
A single training run on MPF.2021.2.8 produces the universal
M3GNet-EFS model.
\lstinputlisting[language=turtle]{artifacts/kg/listings/chen2022/q09-run-and-model.ttl}

\paragraph{Q10 -- Benchmark results.}
Test-set MAEs on the held-out MPF.2021.2.8 split: 35\,meV/atom on
energy, 72\,meV/\AA{} on force components, 0.41\,GPa on stress
components.
\lstinputlisting[language=turtle]{artifacts/kg/listings/chen2022/q10-benchmark-results.ttl}

\paragraph{Q11 -- Computational resources.}
\emph{Not reported} as ontology-encodable data. The paper notes a
single inference example (100-step relaxation of K$_{57}$Se$_{34}$ on
one Intel Xeon E5-2620 v4 core takes ${\sim}22$\,s); recording this as
$\dlRole{inferenceTimePerAtom}$ would require imputing per-atom and
per-step rates, which we leave to Q12.
\lstinputlisting[language=turtle]{artifacts/kg/listings/chen2022/q11-resources.ttl}

\paragraph{Q12 -- Gaps.}
Beyond Q11, the protocol's "one paper $\equiv$ one method $\times$ one
material system" framing is stretched: the universal coverage is
recorded as a single $\MaterialSystemC$ but cannot enumerate the
elements, space groups, or chemistries the model covers without
extending the ontology. DFT settings that vary by MP task (energy
cutoff, k-mesh) are not surfaced. The paper also reports zero-shot
relaxations against a held-out 14{,}055-compound test set and
crystal-structure prediction across selected systems; these
downstream evaluations are not encoded as separate
$\BenchmarkResultC$ instances.
